%In the proposed pipeline, Stage A screens epochs using a union rule over the channels supplied to the pipeline, so changing the channel set can also change which epochs reach Stage B.
%Restricting the input channels therefore changes the detector itself, not merely the signals presented to a fixed detector.
%The EEG channel-subset analysis accounts for this coupling by re-running the complete pipeline separately for each channel subset.
%We first supply the complete 32-channel montage to the pipeline as a baseline, reported in Section~\ref{sec:exp1_reference}. This configuration assumes no prior knowledge of which brain region projects blinks most strongly, making it the most general starting point against which reduced channel sets are compared.
%%
% Section~\ref{sec:exp1_regional}
% then evaluates anatomically defined regional subsets to identify which scalp regions yield the strongest detection performance within the pipeline and to quantify how performance changes when the remaining regions are excluded.
% %
% Section~\ref{sec:exp1_single_channel}
%subsequently assesses individual channels to determine whether a single electrode can retain detection performance comparable to the full montage, which would support a simpler and more practical single-channel system.

In the proposed pipeline, Stage~A screens epochs using a union rule over the channels supplied to the pipeline, so changing the channel set can also change which epochs reach Stage~B. Restricting the input channels therefore changes the detector itself, not merely the signals presented to a fixed detector. The EEG channel-subset analysis accounts for this coupling by re-running the complete pipeline separately for each channel subset. We first supply the complete 32-channel montage to the pipeline as a baseline, reported in Section~\ref{sec:exp1_reference}. This configuration assumes no prior knowledge of which brain region projects blinks most strongly, making it the most general reference for the reduced configurations that follow.
%
Section~\ref{sec:exp1_regional} then evaluates anatomically defined regional subsets, where the union rule operates only over the channels within each region. Comparing these subsets with the full montage identifies which scalp regions yield the strongest detection performance within the pipeline and quantifies how performance changes when the remaining regions are excluded.
%
Section~\ref{sec:exp1_single_channel} takes the reduction further by supplying one electrode at a time, so that the union rule reduces to single-channel screening and both the candidate epochs and the threshold are derived from that electrode alone. Comparison with the full montage then shows whether a single electrode can retain comparable detection performance, which would support a simpler and more practical single-channel system.